\documentclass[11pt]{article}

\usepackage{amsmath}
\usepackage{amsthm}
\usepackage{amssymb}

\newtheorem{theorem}{Theorem}
\newtheorem{lemma}{Lemma}
\newtheorem{proposition}{Proposition}
\theoremstyle{remark}
\newtheorem{remark}{Remark}

\newcommand{\Des}{\operatorname{Des}}
\newcommand{\SYT}{\operatorname{SYT}}
\newcommand{\sort}{\operatorname{sort}}
\newcommand{\ord}{\operatorname{ord}}
\newcommand{\maj}{\operatorname{maj}}

\title{Closing the order law's $\tau=0$ gap: an SYT with all-free descents exists, via a descent-set Kostka number}
\author{Clio}
\date{2026-06-02}

\begin{document}
\maketitle

\begin{abstract}
The order law states $\ord_{x=q^2} Z_\lambda = \min_{T\in\SYT(\lambda)} s(T) = \tau(\tau+1)/2$ for a parity-twisted descent statistic $s$. Its lower bound and the $\tau>0$ achievability were proved earlier; the lone remaining gap was $\tau=0$ achievability --- the existence of a standard Young tableau with no ``paid'' descent. We close it: such a tableau exists for every $\tau=0$ shape, and in fact the number of them is a Kostka number $K_{\lambda,\mu}$, positive exactly when $\tau=0$. The proof is the classical descent-set/Kostka identity plus a two-line dominance computation; we also correct an earlier mis-attribution of this gap to Swanson's modular major-index existence theorem.
\end{abstract}

\section{Setup}

Fix $n\ge 1$ and a partition $\lambda\vdash n$ with $\ell=\lambda'_1$ rows (the length of the first column). For a standard Young tableau $T\in\SYT(\lambda)$, the \emph{descent set} is
\[
\Des(T)=\{\,i\in[1,n-1]: \text{$i+1$ lies in a strictly lower row than $i$ in $T$}\,\}.
\]
Define the \emph{parity-twisted weight}
\[
w_i = \begin{cases} 2i-1, & n-i \text{ odd},\\ 0, & n-i \text{ even},\end{cases}
\qquad
s(T)=\sum_{i\in\Des(T)} w_i .
\]
Call a position $i$ \emph{free} if $n-i$ is even (equivalently $i\equiv n\pmod 2$, so $w_i=0$) and \emph{paid} otherwise. Let
\[
S=\{\,i\in[1,n-1]: i\equiv n\pmod 2\,\}
\]
be the set of free positions. Set
\[
\tau=\max\bigl(0,\;2\ell-n-1\bigr).
\]
Since the paid positions are exactly those carrying a nonzero weight, we have the basic equivalence
\begin{equation}\label{eq:szero}
s(T)=0 \iff \Des(T)\subseteq S \qquad (\text{all descents free}).
\end{equation}

For background (established in the companion note
\texttt{2026-05-31-orderlaw-de\-scent-minimum}): the lower bound $s(T)\ge \tau(\tau+1)/2$ holds for all $T\in\SYT(\lambda)$, and for $\tau>0$ there is an explicit minimizer attaining equality. The remaining gap was the case $\tau=0$: does there exist $T$ with $s(T)=0$? By \eqref{eq:szero} this is exactly the question of whether some $T\in\SYT(\lambda)$ has all of its descents at free positions. We answer it affirmatively and, in fact, count such tableaux.

\section{The descent-set Kostka identity}

\begin{lemma}[classical]\label{lem:kostka}
For any subset $S=\{s_1<\dots<s_k\}\subseteq[1,n-1]$, let
\[
\alpha(S)=(s_1,\;s_2-s_1,\;\dots,\;s_k-s_{k-1},\;n-s_k)
\]
be the composition of $n$ whose set of partial sums (excluding the full sum $n$) is exactly $S$. Then
\[
\#\{\,T\in\SYT(\lambda): \Des(T)\subseteq S\,\} = K_{\lambda,\alpha(S)},
\]
the number of semistandard Young tableaux of shape $\lambda$ and content $\alpha(S)$, equivalently $\langle s_\lambda, h_{\alpha(S)}\rangle$. Since $K_{\lambda,\alpha}$ depends only on the rearrangement of $\alpha$ into a partition, write $\mu=\sort(\alpha(S))$; then the count is $K_{\lambda,\mu}$.
\end{lemma}

\begin{proof}
This is the standard descent-set/Kostka identity; see Stanley, \emph{Enumerative Combinatorics} vol.~2, Cor.~7.19.5 and Thm.~7.19.7, and Gessel. Concretely it is the pairing of the complete-homogeneous expansion $h_\alpha=\sum_\lambda K_{\lambda\alpha}\,s_\lambda$ against the fundamental quasisymmetric expansion $s_\lambda=\sum_{T\in\SYT(\lambda)} F_{\Des(T)}$, together with the fact that the coefficient extracting $\Des(T)\subseteq S$ from $F_{\Des(T)}$ recovers precisely the tableaux whose descent set is contained in $S$.
\end{proof}

\section{The composition for free positions}

\begin{lemma}\label{lem:comp}
For $S=\{\,i\in[1,n-1]: i\equiv n\pmod 2\,\}$:
\begin{itemize}
\item if $n$ is even: $S=\{2,4,\dots,n-2\}$, so $\alpha(S)=(2,2,\dots,2)$ with $n/2$ twos, and $\mu=(2^{n/2})$;
\item if $n$ is odd: $S=\{1,3,\dots,n-2\}$, so $\alpha(S)=(1,2,2,\dots,2)$ (a single $1$ followed by $(n-1)/2$ twos), and $\mu=(2^{(n-1)/2},1)$.
\end{itemize}
In both cases the conjugate partition is
\[
\mu'=\bigl(\lceil n/2\rceil,\ \lfloor n/2\rfloor\bigr),
\]
a partition with at most two parts; in particular $\mu'_1=\lceil n/2\rceil$.
\end{lemma}

\begin{proof}
We read off $\alpha(S)$ from the partial-sum description of Lemma~\ref{lem:kostka}.

\emph{$n$ even.} The free positions are the even integers in $[1,n-1]$, namely $S=\{2,4,\dots,n-2\}$. The consecutive gaps are all $2$, and the final gap is $n-(n-2)=2$, so $\alpha(S)=(2,2,\dots,2)$ with $n/2$ parts. Its partial sums are $2,4,\dots,n$, and intersecting with $[1,n-1]$ recovers $S$, as required. Sorting gives $\mu=(2^{n/2})$.

\emph{$n$ odd.} The free positions are the odd integers in $[1,n-1]$, namely $S=\{1,3,\dots,n-2\}$. The first gap is $1$, the remaining consecutive gaps are $2$, and the final gap is $n-(n-2)=2$, so $\alpha(S)=(1,2,2,\dots,2)$ with a leading $1$ and $(n-1)/2$ twos. Its partial sums are $1,3,5,\dots,n$, whose intersection with $[1,n-1]$ is $S$. Sorting gives $\mu=(2^{(n-1)/2},1)$.

For the conjugate, count column lengths. In either case $\mu$ has parts equal to $2$ and possibly one part equal to $1$. The first column of $\mu$ has one box per part, i.e.\ $\mu'_1$ equals the number of parts: $n/2$ when $n$ is even and $(n-1)/2+1=(n+1)/2$ when $n$ is odd; both equal $\lceil n/2\rceil$. The second column has one box for each part of size $\ge 2$, i.e.\ $\mu'_2$ equals the number of twos: $n/2$ when $n$ is even and $(n-1)/2$ when $n$ is odd; both equal $\lfloor n/2\rfloor$. There is no third column. Hence $\mu'=(\lceil n/2\rceil,\lfloor n/2\rfloor)$.
\end{proof}

\section{Dominance}

We write $\trianglerighteq$ for the dominance order on partitions of $n$.

\begin{lemma}\label{lem:dom}
With $\mu$ as in Lemma~\ref{lem:comp},
\[
K_{\lambda,\mu}>0 \iff \lambda\trianglerighteq\mu \iff \mu'\trianglerighteq\lambda'
\iff \lambda'_1\le\lceil n/2\rceil \iff \tau(\lambda)=0.
\]
\end{lemma}

\begin{proof}
The first equivalence is the Kostka positivity criterion: $K_{\lambda\mu}>0$ if and only if $\lambda$ dominates $\mu$. The second is the standard fact that dominance is reversed under conjugation: $\lambda\trianglerighteq\mu \iff \mu'\trianglerighteq\lambda'$.

For the third equivalence, recall from Lemma~\ref{lem:comp} that $\mu'=(\lceil n/2\rceil,\lfloor n/2\rfloor)$ has only two (possibly nonzero) parts. The dominance condition $\mu'\trianglerighteq\lambda'$ requires
\[
\sum_{i\le j}\mu'_i\ \ge\ \sum_{i\le j}\lambda'_i \qquad\text{for all } j\ge 1 .
\]
For $j=1$ this reads $\lceil n/2\rceil\ge \lambda'_1$. For $j\ge 2$ the left-hand side is already the full sum $\mu'_1+\mu'_2=\lceil n/2\rceil+\lfloor n/2\rfloor=n$, while the right-hand side $\sum_{i\le j}\lambda'_i\le n$; the inequality is therefore automatic. Hence the single binding constraint is $\lambda'_1\le\lceil n/2\rceil$, i.e.\ $\ell\le\lceil n/2\rceil$.

Finally, $\lambda'_1\le\lceil n/2\rceil \iff 2\lambda'_1\le 2\lceil n/2\rceil$. Since $2\lceil n/2\rceil\le n+1$ for both parities (equality when $n$ is odd, and $2\lceil n/2\rceil=n\le n+1$ when $n$ is even), and since $\lambda'_1$ is an integer, $\lambda'_1\le\lceil n/2\rceil$ is equivalent to $2\lambda'_1\le n+1$, i.e.\ to $2\lambda'_1-n-1\le 0$, i.e.\ to $\tau=\max(0,2\lambda'_1-n-1)=0$.
\end{proof}

\section{Main theorem}

\begin{theorem}\label{thm:main}
Let $\lambda\vdash n$ with $\tau(\lambda)=0$, and let $\mu$ be as in Lemma~\ref{lem:comp}. Then
\[
\#\{\,T\in\SYT(\lambda): s(T)=0\,\} = K_{\lambda,\mu}\ \ge\ 1 .
\]
In particular there exists $T\in\SYT(\lambda)$ achieving $s(T)=0=\tau(\tau+1)/2$. Combined with the lower bound $s(T)\ge\tau(\tau+1)/2$ and the explicit $\tau>0$ construction, the order law
\[
\min_{T\in\SYT(\lambda)} s(T)=\tfrac{\tau(\tau+1)}{2}
\]
now holds, with full rigor, uniformly over all partitions $\lambda$.
\end{theorem}

\begin{proof}
By \eqref{eq:szero}, $s(T)=0$ if and only if $\Des(T)\subseteq S$. Applying Lemma~\ref{lem:kostka} with this $S$ counts such tableaux as $K_{\lambda,\alpha(S)}=K_{\lambda,\mu}$, where Lemma~\ref{lem:comp} identifies $\mu$. Since $\tau(\lambda)=0$, Lemma~\ref{lem:dom} gives $K_{\lambda,\mu}>0$, hence the count is at least $1$ and a minimizer with $s(T)=0$ exists. The remaining clause assembles the three legs: the lower bound and $\tau>0$ achievability are established in the companion note, and the present theorem supplies $\tau=0$ achievability, so the minimum equals $\tau(\tau+1)/2$ for every $\lambda$.
\end{proof}

\section{Remarks}

\begin{remark}[Exact count of minimizers]
For $\tau=0$ the number of order-law minimizers is precisely the Kostka number $K_{\lambda,\mu}$ --- a refinement of mere existence. For instance, with $n=7$: the shape $\lambda=(4,2,1)$ gives $K_{\lambda,\mu}=6$, and $\lambda=(5,2)$ gives $K_{\lambda,\mu}=5$.
\end{remark}

\begin{remark}[Conjugate form]
Under transposition one has $s(T')=\binom{n}{2}-s(T)$. Consequently the count above can be read on the conjugate shape as
\[
\#\{\,T'\in\SYT(\lambda'): \text{every paid position is a descent}\,\}=K_{\lambda,\mu}.
\]
\end{remark}

\begin{remark}[Correction of attribution]
This gap was previously conjectured to coincide with Swanson's existence theorem for tableaux with a prescribed modular major index (arXiv:1701.04963). It does not. Swanson controls
\[
\maj(T)=\sum_{i\in\Des(T)} i \pmod{n},
\]
a single linear functional of the descent set, reduced modulo $n$. The condition $s(T)=0$ instead constrains \emph{which positions, by parity}, may be descents. The latter is a finer, positional condition on the descent set that Swanson's mod-$n$ machinery cannot detect: a tableau can realize any target value of $\maj\bmod n$ while still placing a descent at a wrong-parity (paid) position. The correct and far more elementary tool is the descent-set Kostka positivity of Lemma~\ref{lem:kostka}.
\end{remark}

\section{The exact minimizer count and uniform achievability}

The $\tau=0$ argument extends verbatim to all $\tau$ and pins down the exact number of order-law minimizers, replacing the explicit $\tau>0$ construction of the companion note with the same one-line dominance fact.

Write $\ell=\lambda'_1$ (the number of rows) and $K=\lceil\tau/2\rceil$. Recall that the paid positions are those $i$ with $n-i$ odd, the free positions those with $n-i$ even (equivalently $i\equiv n\pmod 2$); let $F$ denote the set of free positions in $[1,n-1]$. Let $p_1<p_2<\cdots$ enumerate the paid positions; the weights $w_{p_1}<w_{p_2}<\cdots$ increase. Define
\[
S^{*} \;=\; F \;\cup\; \{p_1,\dots,p_K\} \qquad\text{(free positions together with the $K$ cheapest paid positions).}
\]

\begin{proposition}[Minimizers are exactly the $\Des\subseteq S^{*}$ tableaux]
For $T\in\SYT(\lambda)$, $\;s(T)=\tau(\tau+1)/2\iff \Des(T)\subseteq S^{*}$.
\end{proposition}
\begin{proof}
($\Leftarrow$) If $\Des(T)\subseteq S^{*}$ then every paid descent lies in $\{p_1,\dots,p_K\}$, so $T$ has at most $K$ paid descents; the lower-bound lemma (Lemma B of the companion note) forces at least $K$. Hence $T$ has exactly $K$ paid descents, necessarily at $p_1,\dots,p_K$, and $s(T)=\sum_{t=1}^K w_{p_t}=\tau(\tau+1)/2$.

($\Rightarrow$) Let $P=\{i_1<i_2<\cdots\}$ be the paid descents of $T$. The lower-bound chain
\[
s(T)=\sum_{t\le|P|}w_{i_t}\ \ge\ \sum_{t\le|P|}w_{p_t}\ \ge\ \sum_{t\le K}w_{p_t}=\tfrac{\tau(\tau+1)}{2}
\]
consists entirely of equalities when $s(T)$ is minimal. Equality in the second step forces $|P|=K$; equality in the first forces $i_t=p_t$ for $t\le K$. So the paid descents are exactly $\{p_1,\dots,p_K\}\subseteq S^{*}$, and the free descents lie in $F\subseteq S^{*}$; thus $\Des(T)\subseteq S^{*}$.
\end{proof}

\begin{theorem}[Exact minimizer count; uniform achievability]
Let $M(\lambda)=\#\{T\in\SYT(\lambda):s(T)=\tau(\tau+1)/2\}$. Then $M(\lambda)=K_{\lambda,\nu}$, where for $\tau>0$
\[
\nu=(2^{\,n-\ell},\,1^{\,\tau+1}),\qquad \nu'=(\ell,\;n-\ell),
\]
and for $\tau=0$, $\nu=\mu$ of the previous sections. In every case $\lambda\trianglerighteq\nu$, so $M(\lambda)\ge1$: the order-law minimum is achieved for all $\lambda$, with no explicit construction.
\end{theorem}
\begin{proof}
By the Proposition and Lemma~\ref{lem:kostka} (the descent-set Kostka identity), $M(\lambda)=\#\{T:\Des(T)\subseteq S^{*}\}=K_{\lambda,\alpha(S^{*})}$. We compute $\alpha(S^{*})$ for $\tau>0$. Since $\tau+1=2\ell-n\equiv n\pmod 2$, the position $\tau+1$ is free. One checks (in both parities of $n$) that the $K=\lceil\tau/2\rceil$ cheapest paid positions are exactly the paid positions $\le\tau$, so
\[
S^{*} \;=\; \{1,2,\dots,\tau+1\}\;\cup\;\{\text{free positions in }[\tau+2,\,n-1]\},
\]
and the free positions in $[\tau+2,n-1]$ are $\{\tau+3,\tau+5,\dots,n-2\}$. Thus the sorted $S^{*}$ is $1,2,\dots,\tau+1,\;\tau+3,\tau+5,\dots,n-2$, whose consecutive gaps give the composition $\alpha(S^{*})=(1^{\,\tau+1},2^{\,n-\ell})$: the initial run $1,\dots,\tau+1$ contributes $\tau+1$ parts equal to $1$, and the subsequent step-$2$ jumps together with the final part $n-(n-2)=2$ contribute $\tfrac{n-1-\tau}{2}=n-\ell$ parts equal to $2$. Rearranged, this is $\nu=(2^{\,n-\ell},1^{\,\tau+1})$, of size $(\tau+1)+2(n-\ell)=n$ (using $\tau+1=2\ell-n$). Its conjugate is $\nu'=(\ell,\,n-\ell)$, a two-part partition. Hence $\lambda\trianglerighteq\nu\iff\nu'\trianglerighteq\lambda'$, whose only binding inequality is $\nu'_1=\ell\ge\lambda'_1=\ell$ --- an equality, hence always satisfied (the $j\ge2$ inequalities read $n\ge\sum_{i\le j}\lambda'_i$). Therefore $K_{\lambda,\nu}\ge1$. For $\tau=0$ the count and the dominance reduce to the previous sections ($\nu'_1=\lceil n/2\rceil\ge\ell\iff\tau=0$). In all cases $M(\lambda)=K_{\lambda,\nu}\ge1$, giving a construction-free, uniform proof of achievability that subsumes the explicit $\tau>0$ construction of the companion note.
\end{proof}

\begin{remark}
The content $\nu$ is always of two-column shape (parts $\le 2$): $\nu=(2^{n-\ell},1^{\tau+1})$ for $\tau>0$ and $(2^{\lfloor n/2\rfloor},1^{\,n-2\lfloor n/2\rfloor})$ for $\tau=0$. The number of order-law minimizers is therefore a Kostka number counting SSYT of two-column content --- e.g.\ $n=7$: $\lambda=(4,2,1)$ gives $M=6$, $\lambda=(5,2)$ gives $M=5$ (both $\tau=0$); for $\tau>0$ shapes $M$ is typically $1$. There is \emph{no} conjugation symmetry $M(\lambda)=M(\lambda')$, since $\tau(\lambda)\ne\tau(\lambda')$ in general.
\end{remark}

\section{Computational verification}

The minimizer-count identity $M(\lambda)=K_{\lambda,\nu}$ with $\nu=(2^{n-\ell},1^{\tau+1})$ was verified with $0$ mismatches over all $194$ partitions $n\le 11$ via independent SSYT enumeration (\texttt{verify\_kostka\_bridge.py} / \texttt{min\_mult\_task2.py}); the equivalence ``minimizers $\iff\Des\subseteq S^{*}$'' was checked with $0/1979$ exceptions. All three legs were verified by an independent script \texttt{verify\_kostka\_bridge.py}, which enumerates semistandard tableaux for the Kostka side using no shared code with the $\SYT$/$s(T)$ side:
\begin{enumerate}
\item[(i)] $\#\{T:s(T)=0\}=K_{\lambda,\mu}$: \textbf{0 mismatches} over all $194$ partitions with $n\le 11$.
\item[(ii)] $\lambda\trianglerighteq\mu \iff \tau=0$: \textbf{0 failures} over all $507$ partitions with $n\le 14$.
\item[(iii)] the single binding constraint $\ell\le\lceil n/2\rceil$: \textbf{0 failures} for $n\le 12$.
\end{enumerate}
Earlier, all $\tau=0$ shapes with $n\le 12$ were brute-confirmed to admit an explicit $s(T)=0$ witness, in agreement with the existence statement of Theorem~\ref{thm:main}.

\end{document}
